%!TeX root=../main.tex
% Figura 2.3: arquitetura conceitual do kworkflow.
% Adaptada de Tadokoro, Siqueira e Meirelles (2025), Fig. 6 ("kw conceptual
% architecture"). Mantém as camadas do original e destaca as duas
% bibliotecas que gravam os dados de uso (kwlib, onde fica statistics_manager,
% e kw_db), retomadas na
% seção 2.4. Como no original, as ligações são ilustrativas.
\begin{tikzpicture}[
    x=1cm, y=1cm,
    every node/.style={font=\footnotesize, align=center},
    peca/.style={draw, rounded corners=3pt, minimum width=1.9cm,
                 minimum height=0.75cm, inner sep=2pt, fill=white,
                 font=\footnotesize\ttfamily},
    dado/.style={peca, very thick, fill=black!12},
    camada/.style={draw, dashed, rounded corners=6pt, inner sep=6pt},
    rotulo/.style={font=\scriptsize\scshape, anchor=west},
    seta/.style={-{Latex[length=1.6mm]}, semithick},
  ]
  % Hub
  \node[peca, minimum width=2.6cm, fill=black!6] (hub) at (4.5, 3.6)
    {kw\\[-1pt]{\normalfont\scriptsize(ponto de entrada)}};

  % Componentes
  \node[peca] (c1) at (0.8, 2.0) {build};
  \node[peca] (c2) at (3.0, 2.0) {deploy};
  \node[peca] (c3) at (5.2, 2.0) {send-patch};
  \node[peca] (c4) at (7.4, 2.0) {report};
  \node (c5) at (9.0, 2.0) {\dots};

  % Bibliotecas
  \node[dado] (l1) at (0.8, 0.4) {kw\_db};
  \node[dado] (l2) at (3.0, 0.4) {kwlib};
  \node[peca] (l3) at (5.2, 0.4) {kw\_string};
  \node[peca] (l4) at (7.4, 0.4) {remote};
  \node (l5) at (9.0, 0.4) {\dots};

  % Plugins
  \node[peca, minimum width=2.7cm] (p1) at (2.6, -1.2) {kernel\_install};
  \node[peca, minimum width=2.3cm] (p2) at (5.4, -1.2) {subsystems};
  \node (p3) at (7.2, -1.2) {\dots};

  % Molduras das camadas
  \node[camada, fit=(c1)(c5)] (fc) {};
  \node[camada, fit=(l1)(l5)] (fl) {};
  \node[camada, fit=(p1)(p3)] (fp) {};
  \node[rotulo] at ([xshift=2pt, yshift=7pt]fc.north west) {componentes};
  \node[rotulo] at ([xshift=2pt, yshift=7pt]fl.north west) {bibliotecas};
  \node[rotulo] at ([xshift=2pt, yshift=7pt]fp.north west) {plugins};

  % O hub despacha para os componentes
  \foreach \c in {c1, c2, c3, c4} {
    \draw[seta] (hub.south) -- (\c.north);
  }
  % Os componentes usam bibliotecas e plugins (ligações ilustrativas)
  \draw[seta] (fc.south) -- (fl.north);
  \draw[seta] (fl.south) -- (fp.north);
\end{tikzpicture}
